\section{Finite certificates and stable exact memory}\label{sec:r73-finite-certificates}

The normal form gives exact realizations at reversible responses. We next prove
that their memory obstructions admit finite tests and persist under a specified
class of physical perturbations. Finite-word equivalence for linear machines is
classical \cite{Tzeng92}; the finite bound below is its direct difference-space
argument. The additional conclusions concern the positive separation from
\emph{all} smaller stochastic machines, its effective algebraic certificate,
and a matching realization under damping toward a common background response.

\subsection{Finite tests against a fixed memory budget}
Suppose a response $b:A^*\to\Delta(C)$ has a supplied real linear representation
\[
 b(w)=\beta T_w B,
 \qquad \beta\in\R^{1\times d},\quad T_a\in\R^{d\times d},
 \quad B\in\R^{d\times |C|}.
\]
Minimality of this representation is not required. For $M\ge1$ let
$\mathcal P_M$ be the compact product of the row-stochastic matrix simplexes
for $(Q_a)_{a\in A}$ and $D\in\R^{M\times |C|}$. Its initial row is the
prescribed pure row $e_1^{\mathsf T}$, as before. Set
\begin{equation}\label{eq:r73-finite-gap}
 L=d+M-1,\qquad
 \delta_M(b)=\min_{(Q,D)\in\mathcal P_M}
 \max_{\substack{|w|\le L\\c\in C}}
 \left|e_1^{\mathsf T}Q_wD_c-b(w,c)\right|.
\end{equation}

\begin{theorem}[Finite obstruction and effective separation]\label{thm:r73-finite-tests}
An $M$-label stochastic machine realizes $b$ on every word if and only if it
does so on the words of length at most $d+M-1$. Consequently
$\delta_M(b)>0$ if and only if no exact $M$-label realization exists.
If another supplied response $\widetilde b$ satisfies
\[
 \max_{|w|\le L,\ c\in C}|\widetilde b(w,c)-b(w,c)|\le\eta,
\]
then every $M$-label machine has error at least $\delta_M(b)-\eta$ on these
same probes. This last assertion imposes no finite-rank hypothesis on
$\widetilde b$.

If the given entries of $\beta,T_a,B$ are real algebraic numbers,
$\delta_M(b)$ is an effectively computable real algebraic number, and a
minimizing machine with real algebraic entries can be computed. When it is
positive, a positive rational number strictly below it and a real-algebraic
quantifier-elimination certificate of that strict lower bound can be computed.
No polynomial-time or practical-size bound is asserted.
\end{theorem}
\begin{proof}
For a fixed candidate machine use the block-diagonal matrices
$S_a=\operatorname{diag}(T_a,Q_a)$, the row
$u=(\beta,-e_1^{\mathsf T})$, and the output columns
$v_c=(B_c,D_c)^{\mathsf T}$. The difference on $w$ is $uS_wv_c$.
Let $R_j$ be the span of $uS_w$ for $|w|\le j$. These spaces increase;
if $R_j=R_{j+1}$, the common space is invariant under every $S_a$ and all
later spaces coincide. Since $u\ne0$, the sequence starts in dimension one
inside a space of dimension $d+M$ and must have stabilized by $j=d+M-1$.
Vanishing on the indicated finite set therefore implies vanishing on all
words. This is the same linear-equivalence principle as the inherited
finite-realization test.

The maximum in \eqref{eq:r73-finite-gap} is a continuous function on the
compact set $\mathcal P_M$, so its minimum is attained. It is zero precisely
when one machine passes every finite probe, hence, by the first assertion,
realizes every word. The perturbation claim follows pointwise from the triangle
inequality and then from the minimum defining $\delta_M(b)$.

For algebraic input, every probe value is algebraic. The constraints defining
row-stochasticity and the inequalities between the finitely many polynomial
errors and a variable $t$ are polynomial constraints over the real algebraic
numbers. Real quantifier elimination decides their feasibility and computes
the least feasible $t$ together with an algebraic witness. A singleton
semialgebraic set over the real algebraic numbers consists of an algebraic
number. If $\delta_M(b)>0$, effective isolation supplies a positive rational
$\sigma<\delta_M(b)$; infeasibility with all errors at most $\sigma$ is a
finite real-algebraic certificate. This uses the classical decision procedure
\cite{Seidenberg54}; it does not assert that this procedure was executed by
the finite diagnostics accompanying this article.
\end{proof}

For the prepared quantum response, these finite probes have an explicit
physical perturbation bound. Suppose two density matrices differ in trace
norm by at most $h_\rho$, corresponding channels differ in diamond norm by
at most $h_\Phi$, and corresponding terminal effects differ in operator norm
by at most $h_E$. All channels are trace preserving and completely positive,
and both lists of effects are valid measurements. Contractivity and telescoping
give $\|\widetilde\Phi_w(\widetilde\rho)-\Phi_w(\rho)\|_1
\le h_\rho+|w|h_\Phi$. The state difference has trace zero; subtracting
$I/2$ from a terminal effect, whose spectrum lies in $[0,1]$, gives
\begin{equation}\label{eq:r73-physical-probes}
 \max_{|w|\le L,c}|\widetilde b(w,c)-b(w,c)|
 \le h_E+\tfrac12(h_\rho+Lh_\Phi).
\end{equation}
Thus a strict smaller-memory obstruction survives a specified open
neighbourhood of physical primitives. This supplies a lower bound there;
it does not assert that an arbitrary perturbation has a realization of the
same size.

\subsection{A matching upper under a common background}
Let $b$ be compactly reversible, nonconstant, and have finite stochastic
memory $m=m(b)$. Let $G$ be its finite observable group, with the chronological
convention of the normal-form theorem, and put
\[
 \overline b=|G|^{-1}\sum_{g\in G} b(g),\qquad
 C_b=\max_{g\in G,c\in C}|b(g,c)-\overline b(c)|>0.
\]
Choose a damping factor $\alpha_a\in[0,1]$ for each command, put
$\beta_w=\prod_{j=1}^{|w|}\alpha_{w_j}$, with $\beta_\varnothing=1$, and define
\begin{equation}\label{eq:r73-common-background}
 b_{\boldsymbol\alpha}(w)
   =\beta_w b(w)+(1-\beta_w)\overline b.
\end{equation}

\begin{theorem}[Stable exact memory under background damping]\label{thm:r73-memory-plateau}
The response \eqref{eq:r73-common-background} has an exact realization with
$m$ labels for every choice of the damping factors. Put $L=d+m-2$, where
$d$ is any supplied representation dimension of $b$, and let
$\epsilon=\max_a(1-\alpha_a)$. Then
\begin{equation}\label{eq:r73-memory-plateau}
 \epsilon<\min\left\{1,\frac{\delta_{m-1}(b)}{LC_b}\right\}
 \quad\Longrightarrow\quad m(b_{\boldsymbol\alpha})=m.
\end{equation}
For algebraic input satisfying the hypotheses, a positive algebraic cutoff,
or a smaller positive rational cutoff, is effectively computable.
For a constant response every damping leaves that response unchanged and its
minimum is one label; the expression involving $m-1$ is not used.
\end{theorem}
\begin{proof}
Take a minimal realization in the normal form proved above. On its recurrent
labels each $Q_a$ is a permutation, and an extra initializer is present only
when the initial distribution on these labels is mixed. Average that
distribution over $G$, and call the resulting row $\pi$, extended by zero
on the initializer when present. It obeys
\[
 \pi Q_a=\pi,\qquad \pi D=\overline b.
\]
Indeed, the group average is invariant under each permutation. The initializer
has no incoming transition, and its first-input row is the original initial
distribution followed by the commanded permutation, so the same identities
hold on the entire machine. Set $J=\mathbf1\pi$. Then
$J^2=J$, $Q_aJ=JQ_a=J$, and
\[
 Q_{a,\boldsymbol\alpha}=\alpha_aQ_a+(1-\alpha_a)J
\]
is row-stochastic on the same labels. Multiplication gives, for every nonempty
word and also trivially at the empty word,
\[
 Q_{w,\boldsymbol\alpha}=\beta_wQ_w+(1-\beta_w)J.
\]
Reading with the same decoder proves the upper.

By minimality $\delta_{m-1}(b)>0$. For $|w|\le L$, the elementary product
bound $1-\beta_w\le |w|\epsilon$ gives an error from $b$ of at most
$LC_b\epsilon$. Theorem~\ref{thm:r73-finite-tests} rules out every
$(m-1)$-label realization when this is smaller than $\delta_{m-1}(b)$.
Machines with fewer labels can be padded by unreachable rows, so they too are
excluded. This and the upper prove equality. A nonconstant response has
$m\ge2$ and $d\ge2$, hence $L>0$. The effective assertion follows from the
finite group enumeration and algebraic synthesis theorem, together with
Theorem~\ref{thm:r73-finite-tests}. Finite maxima and averages over the
computed group preserve algebraicity.
\end{proof}

There is a physical realization of the common background for prepared unitary
commands. Let $\mathcal K$ be the compact closure of their physical conjugation
group. The compact convex hull of $\{\Phi(\rho):\Phi\in\mathcal K\}$ has a
unique element $\rho_*$ of least Hilbert--Schmidt norm: strict convexity gives
uniqueness. Every conjugation in $\mathcal K$ preserves that set and norm, so
$\rho_*$ is invariant. Its observable image belongs to the convex hull of the
finite orbit of $b$ and is invariant. Averaging any element of that orbit
over $G$ gives $\overline b$; applying this averaging projection to an
invariant element in its convex hull shows that its observable image is
exactly $\overline b$. Therefore the physical channels
\[
 \widetilde\Phi_a(X)=\alpha_aU_aXU_a^*
                   +(1-\alpha_a)\operatorname{tr}(X)\rho_*
\]
produce \eqref{eq:r73-common-background} under every command word. In general
$\rho_*$ need not be maximally mixed. The theorem concerns damping toward
this common invariant background, not every depolarizing or noisy channel.

\subsection{Periodic plateaux and an explicit rational certificate}
For the periodic mixture of Section~\ref{sec:r72-periodic}, let the distinct
nonreal frequencies have positive amplitudes $\gamma_j$, total
$\Gamma=\sum_j\gamma_j$, and let $M_*$ be its exact partition-and-anchor
minimum. Its dimension is $d=2r+1$ and $C_b\le\Gamma/2$. Hence a common
damping $\alpha=1-\epsilon$ retains exactly $M_*$ labels whenever
\begin{equation}\label{eq:r73-periodic-cutoff}
 0\le\epsilon<
 \min\left\{1,\frac{2\delta_{M_*-1}(b)}{\Gamma(2r+M_*-1)}\right\}.
\end{equation}
This is an effective finite certificate for each algebraic periodic input.
It closes the previously existential plateau cutoff for this class, without
claiming a sharp transition point or a fixed-memory risk asymptotic.

For clarity, one certificate can be calculated without a general
quantifier-elimination run. Consider the unary response
\begin{equation}\label{eq:r73-period-four}
 b_\alpha(n)=\tfrac12\{1+\alpha^n\cos(\pi n/2)\},\qquad n\ge0.
\end{equation}

\begin{proposition}[A certified four-label interval]\label{prop:r73-period-four}
For the undamped sequence $b=b_1$ and the dimension $d=3$,
$\delta_3(b)>1/36$. In particular the exact stochastic memory of
\eqref{eq:r73-period-four} is four for $89/90\le\alpha\le1$.
This interval is a certified sufficient interval, not a claimed sharp one.
\end{proposition}
\begin{proof}
Let $Q$ be any row-stochastic $3\times3$ matrix, put $t=\operatorname{tr}Q$
and $h=\det Q$, and write $y_n=e_1^{\mathsf T}Q^nD$ for any scalar decoder.
Since $Q\mathbf1=\mathbf1$, its characteristic polynomial is
\[
 z^3-tz^2+(t+h-1)z-h.
\]
Thus $y_{n+3}-ty_{n+2}+(t+h-1)y_{n+1}-hy_n=0$. We have $0\le t\le3$,
$|h|\le1$, and $|t+h-1|\le3$. The determinant bound follows directly by
bounding its absolute permutation sum by the product of the row sums.
If the first five values of $y$ are within $\eta$ of
$(1,1/2,0,1/2,1)$, substitution at $n=0,1$ yields
\[
 |t-h|\le16\eta,\qquad |2-t-h|\le16\eta,
 \qquad |t-1|,|h-1|\le16\eta.
\]
Nonnegativity and the row sums imply a further exact identity. For
$a_i=1-Q_{ii}$ it is
\begin{equation}\label{eq:r73-three-state-certificate}
 t^2-3t+3-3h
 =\tfrac12\sum_{i<j}(a_i-a_j)^2
       +3\sum_{i<j}Q_{ij}Q_{ji}\ge0.
\end{equation}
To verify it, the second characteristic coefficient is both $t+h-1$ and
$\sum_{i<j}(Q_{ii}Q_{jj}-Q_{ij}Q_{ji})$; substitute
$Q_{ii}=1-a_i$ and simplify. On the other hand, putting $u=t-1$ and $v=h-1$
gives
\[
 t^2-3t+3-3h=-2-u+u^2-3v
       \le-2+64\eta+256\eta^2.
\]
For $\eta\le1/36$ the right side is at most $-2/81<0$, a contradiction.
The compact minimum on the first five probes is therefore strictly greater
than $1/36$, and so is \eqref{eq:r73-finite-gap}, which here uses the six
probes $0\le n\le5$.

The undamped response has dimension three and a four-cycle realization with
decoder $(1,1/2,0,1/2)$. The same polynomial certificate excludes three labels,
so its minimum is four. Damping that cycle toward the uniform stationary row
realizes \eqref{eq:r73-period-four} with four labels. Its difference from the
undamped sequence on $0\le n\le5$ is at most
$\tfrac52(1-\alpha)$. For $\alpha\ge89/90$ this is at most $1/36$, strictly
below $\delta_3(b)$. The finite obstruction proves the lower and completes
the proof.
\end{proof}

The certificate uses a stochastic-matrix inequality rather than a sample
search over three-label machines. Its accompanying rational diagnostics check
the displayed polynomial identity, the constructed all-word damping
identities on explicit finite models, and exact finite recurrence consequences.
They do not replace the proofs, compute all algebraic gaps, or give an
independent human or editorial assessment.
